% Provenance: copied from the design workspace CambridgePQ/recursion.fast, `report/sections/10-cost.tex`, on 2026-10-19. Read-only copy; the source is authoritative.
% 10-cost -- written by W3 (revised for cost model v0.3)
\section{Cost function and incentives}\label{sec:cost}

The score must track what recursion pays, and it must be a kernel-checked theorem, not a
measurement. Earlier golf competitions priced every RISC-V instruction alike, and agents optimised
that model rather than the problem. This section gives cost model v0.3 (DESIGN §15), which replaces
the flat per-instruction weights of v0.2. It proceeds in order:
\begin{itemize}[nosep]
\item the evidence that a flat per-instruction model is overfit;
\item the two machines the competition prices: a zkVM for the recursion step and a multicore CPU for
  native verification;
\item the one proved object: an 18-counter cost vector over accepting runs;
\item the two boards built on it: the recursion board, a geometric mean over four calibrated zkVM
  profiles, and the latency board, a deterministic timing model with fork-join harts;
\item the parts of v0.2 that carry over: the accepting-run bound, honest preprocessing, caps, the
  Fiat--Shamir surcharge and the blow-up factor;
\item what it costs to implement, and what remains unrealistic.
\end{itemize}
Fields are $K=\mathrm{GF}(2^{64})$ and $E=K[y]/(y^3+y+1)=\mathrm{GF}(2^{192})$
(\cref{subsec:bg-binary}). The machine is RV64IM plus the carry-less-multiply extension Zbc
(\cref{subsec:proto-machine}). Kernel costs come from one emulator trace priced under nine models at once.
Trajectory numbers are model estimates, accurate to about a factor of 2.

%-----------------------------------------------------------------------------------------------
\subsection{Diagnosis: how a flat per-instruction model is overfit}\label{subsec:cost-diagnosis}

\paragraph{What the golf competitions produced.} sig.golf scores one cycle per instruction, four per
multiplication, and treats code as free up to 1\,MiB. Its accepted verify image has 145{,}952
instructions.
\begin{itemize}[nosep]
\item 98{,}134 of them are NOPs (\textbf{67\%}): padding that keeps code at fixed 256-word pitches so
  that table-driven dispatch can \texttt{jalr} into slots.
\item Of the 47{,}818 remaining words, \textbf{59\% are stores and 28\% are \texttt{jal}s}. The code is
  straight-line and table-driven: it stages constants into memory and hops between slots. It contains
  867 loads, 345 ALU operations and 279 branches in all.
\end{itemize}
That code is optimal only where a store and a jump each cost 1 and code size, instruction cache and
paging cost nothing. On RISC~Zero, 584\,KB of straight-line code would page in at about 1{,}130 cycles
per KiB, i.e.\ about 660k cycles against a claimed 11{,}687.

ots.golf counts \textbf{one hash compression as one instruction}. A full-length acceptance is 192
compressions plus 149 instructions, so hashing is 56\% of the score. At a software BLAKE2s cost of
1{,}560--2{,}770 cycles it would be more than 99.9\% of the time. One ots.golf track's record moved from
49{,}335 to 1{,}088 within a week: the agents converged on the model, not on the problem.

\paragraph{Mispricings on this competition's kernels.} One trace of each kernel was priced under nine
models:
\begin{itemize}[nosep]
\item \texttt{uni}, the v0.2 weights;
\item five zkVM weightings: \texttt{zisk}, \texttt{sp1} (no Zbc chip, as SP1 is today), \texttt{sp1z}
  (with a hypothetical Zbc chip), \texttt{r0} (RISC~Zero) and \texttt{jolt};
\item a generic memory-priced zkVM, \texttt{zkmem};
\item an in-order dual-issue scoreboard with caches (\texttt{inord}), using SiFive-7 and SpacemiT-X60
  latencies from GCC's scheduler models;
\item an idealised 4-wide out-of-order dataflow model (\texttt{ooo}).
\end{itemize}
The v0.2 model gets five things wrong:
\begin{itemize}[nosep]
\item \emph{Multiplication.} MUL $=$ CLMUL $=4$ is about 3 times the area ratio of real zkVMs (ZisK
  1.29, SP1 1.38). Field kernels cost 0.7--0.8 times v0.2 under every Zbc-equipped zkVM profile, and
  0.5 times on an in-order core.
\item \emph{Memory and paging.} Loads cost the same as ALU operations, and first touches and paging
  cost nothing. A random 8-byte gather over 4\,MiB costs 13 under \texttt{uni} and 87 on the in-order
  core. On that core the trivial baseline at $2^{20}$ costs 3.2 times its v0.2 count, 65\% of it DRAM
  stalls. One MiB of proof read once costs about 1.16M cycles on RISC~Zero, 9 times the v0.2 byte
  charge.
\item \emph{Code size.} Free code means unrolling wins under \texttt{uni} (1.66 times cheaper) and loses
  under every realistic model: 6 times dearer in-order, 17 times out-of-order, 1.4 times on RISC~Zero.
\item \emph{Hashing.} The virtual Merkle level at 43 is 2.5--5 times below zkVM precompile prices and 37
  times below a native core.
\item \emph{No instruction-level parallelism.} Independent multiplication chains run 1.5--1.9 times faster
  on a real out-of-order core, and no per-instruction count sees it.
\end{itemize}
On field kernels every realistic model is a near-uniform 0.5--0.8 of v0.2. On memory-, control- and
code-bound kernels the ratios are 1.2--8. The flat model's \emph{relative} error across kernel
classes therefore spans about \textbf{15 times}. Seven of thirteen paired kernels change winner under
at least one model (\cref{tab:cost-flips}). Three of those flips are robust to the pipeline parameters:
branchless code under \texttt{ooo}, code size, and tables without Zbc.

\begin{table}[t]
\centering\footnotesize
\setlength{\tabcolsep}{3.5pt}
\begin{tabular}{@{}>{\raggedright\arraybackslash}p{4.6cm}rrrrrrrr@{}}
\toprule
Pair (cost of A $/$ cost of B; $>1$: B wins) & uni & zisk & r0 & jolt & sp1 & sp1z & inord & ooo \\
\midrule
$E$ mult.: Karatsuba / schoolbook & 0.89 & 1.00 & 1.03 & 1.00 & 0.68 & 0.99 & 0.88 & 1.00 \\
$K$ inverse: Fermat / Itoh--Tsujii & 1.71 & 1.70 & 1.50 & 1.69 & 1.76 & 1.70 & 1.49 & \textbf{0.78} \\
$K$ square: clmul / 16\,KiB tables & 0.56 & 0.51 & 0.55 & 0.45 & \textbf{6.14} & 0.49 & 0.54 & 0.79 \\
dispatch: jump table / if-chain & 0.67 & 0.63 & 0.74 & 0.77 & 0.68 & 0.68 & \textbf{1.11} & \textbf{1.18} \\
branchy / branchless (random bit) & 0.66 & 0.63 & 0.77 & 0.69 & 0.64 & 0.64 & 0.86 & \textbf{2.56} \\
staging: unrolled (24\,KiB code) / loop & 0.60 & 0.53 & \textbf{1.37} & 0.60 & 0.62 & 0.62 & \textbf{5.88} & \textbf{16.7} \\
loads: random 4\,MiB / sequential & 3.25 & 3.58 & 6.40 & 3.25 & 3.42 & 3.42 & 21.6 & 11.7 \\
\bottomrule
\end{tabular}
\caption{Paired kernels under nine cost models (selection; \texttt{zkmem} omitted). Bold entries change
the winner relative to \texttt{uni}; the last row keeps its winner but changes its magnitude 7-fold
[D4 §1.2; probe \texttt{d4-probes}]. The jump-table flip is sensitive to the pipeline parameters; the
branchless, code-size and no-Zbc flips hold under every setting.}
\label{tab:cost-flips}
\end{table}

\paragraph{Architecture or engineering?} Re-running the trajectory model under the nine models
(\cref{subsec:bench-trajectory}, \cref{tab:bench-models}) shows that v0.2 ranks \emph{architectures}
tolerably: the top three designs are stable. It mis-ranks \emph{engineering} and mis-states
\emph{where the cost is} (hashing on native cores, bytes on RISC~Zero). The fix is cheap for the zkVM
board, by re-weighting, and needs a timing model for the CPU board.

%-----------------------------------------------------------------------------------------------
\subsection{Candidate cost models}\label{subsec:cost-candidates}

\Cref{tab:cost-candidates} lists the candidates. Four requirements decide between them.
\begin{itemize}[nosep]
\item \emph{Dimension.} The score must be monotone in what users pay.
\item \emph{Ungameability.} No prover-side or benchmark-side choice may move it.
\item \emph{Realism.} It must resist the overfitting described above.
\item \emph{Provability.} It must follow from a theorem whose proof technique exists.
\end{itemize}

\begingroup\small
\renewcommand{\arraystretch}{1.12}
\begin{longtable}{@{}>{\raggedright\arraybackslash}p{0.27\textwidth}>{\raggedright\arraybackslash}p{0.22\textwidth}>{\raggedright\arraybackslash}p{0.45\textwidth}@{}}
\caption{Candidate cost models and verdicts.}\label{tab:cost-candidates}\\
\toprule
Candidate & Form & Verdict \\
\midrule
\endfirsthead
\toprule
Candidate & Form & Verdict \\
\midrule
\endhead
Brief: performance $\div$ size & cycles $/$ bytes & rejected: dimensionally inverted, and any size term rewards obfuscation. Replaced by a hard 1\,MiB cap \\
sig.golf product & bytes $\times$ cycles & rejected for the primary board: fixes an arbitrary exchange rate. Kept as the shape of the proof-size board \\
Flat per-instruction weights (v0.2) & $(1;\ \mathtt{MUL}=\mathtt{CLMUL}=4)$ & superseded: overfit, with a relative error spanning about 15 times across kernel classes (\cref{subsec:cost-diagnosis}). Kept as an unscored display column \\
One calibrated zkVM profile & $w_P\cdot c$ & rejected alone: invites overfitting to one zkVM, whose costs also change between releases \\
\textbf{Four calibrated profiles, geometric mean} & $\bigl(\prod_P\prod_kS_{P,k}\bigr)^{1/12}$ & \textbf{adopted} for the recursion board (\cref{subsec:cost-zk}) \\
Minimum or maximum over profiles & $\min_P$, $\max_P$ & $\min$ rewards overfitting. $\max$ (minimax) is defensible but brittle to one mis-calibrated profile \\
\textbf{Deterministic in-order timing model + fork-join harts} & $T_{16}=W/16+S$ & \textbf{adopted} for the latency board (\cref{subsec:cost-lat}) \\
Cache and DRAM model in the proof & first-touch line sets & deferred: provable only for statically known address ranges; first-touch \emph{counts} are adopted instead \\
Wall time on reference hardware & seconds & calibration and display only: noise of $\pm2$--$5\%$ against sub-percent record margins, not a theorem, and it would put candidate code on the score's trust path \\
Wall time of a 2$\to$1 recursion step & seconds & deferred: needs a trusted arithmetiser (T3, \cref{subsec:rec-t3}) \\
Measured median on hash chains & -- & rejected as a score: gameable by special-casing \\
All-runs or honest-runs bound & see \cref{tab:cost-clauses} & superseded by the accepting-run bound \\
\bottomrule
\end{longtable}
\endgroup

\begin{decision}{COST-1}{Two machines, two boards, one proved vector}
The \emph{recursion board} (machine M-zk) is primary. It prices the prover cost of the verifier's
trace in a zkVM, as a geometric mean over four calibrated zkVM profiles. The \emph{latency board}
(machine M-cpu) is a separate leaderboard, scored from season~1. It prices native multicore
verification, which matters for the decider and for node verification, with a deterministic
pipeline-and-harts model. Both are computed from proved bounds. Wall time is calibration and display
only, and the v0.2 flat weights are a display column. \emph{Rationale:} the two machines reward
different code (branches, tables, unrolling). A combined score would let one hide the other.
\emph{Answers:} user input U3 (flat per-instruction costs are overfit; parallelism is not accounted for). [DESIGN
§15.1; D4-1--D4-3]
\end{decision}

%-----------------------------------------------------------------------------------------------
\subsection{The proved object: an 18-counter vector over accepting runs}\label{subsec:cost-vector}

The cost theorem bounds a vector of monotone counters, not a scalar:
\begin{multline*}
  c=(n_{\mathrm{add}},n_{\mathrm{logic}},n_{\mathrm{shift}},n_{\mathrm{mul}},n_{\mathrm{div}},
     n_{\mathrm{clmul}},n_{\mathrm{ld}},n_{\mathrm{ldsub}},n_{\mathrm{st}},n_{\mathrm{stsub}},
     n_{\mathrm{br}},n_{\mathrm{jmp}},n_{\mathrm{sys}},\\
     m_{\mathrm{words}},m_{\mathrm{pages}},q^{\mathrm{count}},q^{\mathrm{lev}},q^{\mathrm{bytes}}).
\end{multline*}
\begin{itemize}[nosep]
\item Instruction counters separate addition, logic, shifts, multiplication, division and Zbc, and
  split loads and stores into doubleword and sub-doubleword accesses. Branches and jumps are counted
  apart.
\item $m_{\mathrm{words}}$ and $m_{\mathrm{pages}}$ count \emph{first-touched} doublewords and 1\,KiB
  pages. They cover code and data, and \texttt{QUERY} destinations count like stores. In the Lean
  machine, \lean{execute} threads three \lean{Finset}s (doublewords, pages, written pages), and the
  counters are their cardinalities at the end of the run (\lean{Acct.vec}; a written page counts twice).
\item $q^{\mathrm{count}}$, $q^{\mathrm{lev}}=\sum\mathsf{depth}_j$ and
  $q^{\mathrm{bytes}}=\sum\mathsf{leafbytes}_j$ are summed over the \texttt{QUERY}s of a run.
\end{itemize}
Every profile is a non-negative linear function of $c$. A weighted sum of componentwise maxima bounds
the maximum of the weighted sums: $\sum_iw_i\max_{\mathrm{run}}n_i\ge\max_{\mathrm{run}}\sum_iw_in_i$.
So \emph{one} certificate is a valid bound under every profile and under every later re-weighting.
Proving it uses the same loop invariants as before, carried with a tuple. For the first-touch
counters the library has per-touch cardinality lemmas; a footprint combinator for loops that write
$[\mathit{base},\mathit{base}+8n)$ is planned (L-10), and the baseline wrote its own invariant.
The clause is the contract's (\cref{subsec:app-cost} has the whole file):

\begin{leancode}[label={\footnotesize Contract/Cost.lean}]
def Protocol.AcceptingCostAtMost {sz : Sizes} (P : Protocol sz) (img : Image) (c : CostVec) : Prop :=
  ∀ (i : Index sz) x (tr : FullTranscript P.shape) v s, run img (P.tape i x tr) (P.oracles i tr) =
      some (v, s) → v ≤ c
\end{leancode}
Here \lean{run} returns \lean{some} exactly on acceptance, with the 18 counters and the final state;
the tape holds the sizes, the honest \lean{VK i}, $x$ and the challenges of \lean{tr}, and the
oracles are the honest \lean{IDX i} followed by the raw messages of \lean{tr}. So the quantifier
ranges over every index of the size, every public input, every raw oracle content and every
challenge vector, and \lean{≤} is componentwise.

\phantomsection\label{subsec:cost-u2}%
\paragraph{Fuel.} Each hart runs for at most $F=2^{36}\approx6.9\times10^{10}$ steps, and fuel
exhaustion is a rejection. No separate termination theorem is needed: $\mathsf{Agrees}$ already forces
every spec-accepted transcript, including every honest one, to be accepted within fuel. $F$ is a
contract constant that counts executed instructions (steps), not weighted cost. It is per hart: a
child runs with its parent's remaining fuel, and the parent is not charged for the child's steps, so
a run with harts can take up to $(\#\mathrm{spawns}+1)\cdot2^{36}$ steps (A6-F1, row H12). The certified baseline's
accepting runs at $2^{24}$ take $3.09\times10^9$ steps, 4.5\% of $F$ (\cref{sec:baseline}); a succinct
design takes about $10^{6}$--$10^{7}$ steps, so $F$ leaves a margin of $10^4$ or more.

\paragraph{Why accepting runs, and exploits attempted.} \Cref{subsec:obj-what} gives the argument: only
accepting executions are ever arithmetised, the clause covers every honest run and every accepting
run on an adversarial transcript, and a verifier cannot make acceptance cheap selectively without
breaking soundness. \Cref{tab:cost-u2-exploits} lists the attempts made in design note D3; none
gains anything over the all-runs clause, and re-weighting does not touch the argument.

\begingroup\small
\renewcommand{\arraystretch}{1.1}
\begin{longtable}{@{}p{0.03\textwidth}>{\raggedright\arraybackslash}p{0.36\textwidth}>{\raggedright\arraybackslash}p{0.55\textwidth}@{}}
\caption{Attempted exploits of the accepting-run clause [D3 §4.3].}\label{tab:cost-u2-exploits}\\
\toprule
\# & Attempt & Outcome \\
\midrule
\endfirsthead
\toprule
\# & Attempt & Outcome \\
\midrule
\endhead
1 & Quick accept for honest-looking proofs, long check otherwise & every accepting run is bounded; a sound quick path is just a cheaper verifier \\
2 & Optional cheap \enquote{easy mode} protocol chosen by the prover & the maximum over both modes is scored \\
3 & Move work into pre-checks that reject & the accepting path executes every pre-check that passes \\
4 & Loop trip counts read from the proof & bounded on accepting runs by the consistency check; rejecting runs stop at $F$ \\
5 & Data-dependent \texttt{QUERY} counts & $q^{\mathrm{count}}\le Q_k$ on accepting runs \\
6 & Vacuity on unsatisfiable indices & harmless: unprovable statements are never verified, and honest runs accept on every satisfiable index \\
7 & Make accepting runs on false statements expensive & inside the bound; see 1 \\
8 & Trap to escape the bound on a long path & a fault is a reject; acceptance requires \texttt{HALT} with $a_0=1$ \\
9 & T2: cheap output accumulator, defer everything & output shape fixed and pinned by $\mathsf{Agrees}$; deferral is paid at $S_{\mathrm{dec}}/64$; the decider is bounded only by the fuel \\
10 & Grind the demonstration run & demonstrations are unscored; the badge checks measured $\le$ claimed \\
\bottomrule
\end{longtable}
\endgroup

\paragraph{Pitfalls.} (1) \emph{Denial of service.} A verifier may spend up to $F$ steps rejecting,
about 10--20\,s natively, and up to $(\#\mathrm{spawns}+1)\cdot F$ with harts (H12). The planned
mitigations, not in the v0.3 contract, are an optional displayed all-runs badge and, for T2
deciders, a required bound $C^{\mathrm{all}}_{\mathrm{dec}}\le8\,C_{\mathrm{dec}}$ (R-3).
(2) \emph{Proof burden.} After the honest-preprocessing fix, the
accepting and all-runs clauses differ by 5--15\% of the machine-level proof. The main saving is that no
termination theorem is proved [A3 n2]. (3) \emph{Hygiene.} There is one componentwise bound on one
vector, and every score is a Lean function of it (\cref{subsec:cost-zk}).

\begin{table}[t]
\centering\small
\begin{tabularx}{\linewidth}{@{}>{\raggedright\arraybackslash}p{2.6cm}>{\raggedright\arraybackslash}X>{\raggedright\arraybackslash}p{3.4cm}>{\raggedright\arraybackslash}X@{}}
\toprule
Clause & Quantifies over & Precedent & Deployment meaning \\
\midrule
All runs & every input and oracle & ots.golf RISC-V \lean{cycles} & cost bound plus DoS bound; ranking arguments on every data-driven loop \\
\textbf{Accepting runs (adopted)} & every input on which the machine accepts & ots.golf hinted track (every prover-chosen view, accepting executions) & exactly the recursion cost; DoS bounded only by $F$ per hart \\
Honest runs & honest transcripts under every fixed oracle & sig.golf \lean{verificationCycles} & too weak: an accepted dishonest transcript may run longer \\
\bottomrule
\end{tabularx}
\caption{Cost clauses and precedents (\cref{sec:prior}).}
\label{tab:cost-clauses}
\end{table}

\paragraph{Honest preprocessing.} The first design quantified the cost clause over every oracle,
including oracle~0, the index. Under that quantifier the baseline's claimed bound is \emph{false}
[A3-C2]. The baseline reads row and nonzero counts from $\vk$ and $\idx$.
\begin{itemize}[nosep]
\item A huge count makes a loop walk memory until a load faults. Under the all-runs clause that alone
  costs about $10^{10}$ steps in every class.
\item The accepting-run clause does not rescue it. Filling oracle~0 to its admitted length $2^{k+10}$
  with zero-valued entries makes every row hold trivially. The run then \emph{accepts} after about
  $20\cdot2^{k+9}$ steps, about 22 times the claimed $C_k$.
\end{itemize}
The fix loses nothing. $\vk_i$ and $\idx(i)$ are outputs of the deterministic, computable Lean
indexer. The harness computes them and binds $\mathrm{MTroot}(\idx(i))$ and $\vk_i$ into the
Fiat--Shamir prefix, and $\mathsf{Secure}$ and $\mathsf{Agrees}$ are already stated for them. Every
machine clause therefore quantifies over honest preprocessing ($\forall i$ of class shape). The prover
and environment inputs stay adversarial: $x$, every raw oracle content and every challenge. In
deployment this is the standard holographic assumption that the verifier holds the honest key. As a
side effect, trip counts come from $k$ and $\vk_i$, never from prover data, so cost proofs need no
ranking arguments.

%-----------------------------------------------------------------------------------------------
\subsection{The recursion board: four calibrated zkVM profiles}\label{subsec:cost-zk}

\paragraph{Profiles.} Each profile $P$ is a weight vector $w_P$ on the instruction and memory
counters, a hash price $P_{H,P}$ per 64-byte compression, and a byte charge $\lambda_P$. The weights were
normalised from primary cost tables:
\begin{itemize}[nosep]
\item ZisK: area $\times$ operation from \texttt{zisk\_ops\_costs.rs} and the memory state-machine
  columns;
\item SP1 Hypercube: trace area per event from \texttt{rv64im\_costs.json}, with fetch and decode charged
  per executed instruction and a hypothetical Zbc chip priced like MUL;
\item RISC~Zero: one cycle per instruction, 2 for bitwise and division, 1{,}130 cycles per 1\,KiB page
  touched (the documented average of 1{,}094--5{,}130);
\item Jolt: lookup-uniform, with virtual division sequences and expensive sub-word stores.
\end{itemize}
\Cref{tab:cost-profiles} gives the weights.

\begin{table}[t]
\centering\footnotesize
\setlength{\tabcolsep}{4pt}
\begin{tabular}{@{}lrrrrr@{}}
\toprule
Counter & ZK-A ZisK & ZK-B SP1-H & ZK-C RISC~Zero & ZK-D Jolt & v0.2 (display) \\
\midrule
$n_{\mathrm{add}}$ (add[w], sub[w], addi[w], lui, auipc, fence) & 0.73 & 1.00 & 1 & 1 & 1 \\
$n_{\mathrm{logic}}$ (and, or, xor, slt, with imm.) & 1.00 & 1.10 & 2 & 1 & 1 \\
$n_{\mathrm{shift}}$ & 0.97 & 1.27 & 1 & 1 & 1 \\
$n_{\mathrm{mul}}$ & 1.29 & 1.38 & 1 & 1 & 4 \\
$n_{\mathrm{div}}$ & 1.29 & 2.64 & 2 & 8 & 4 \\
$n_{\mathrm{clmul}}$ & 0.97 & 1.38\,[J] & 1\,[J] & 1\,[J] & 4 \\
$n_{\mathrm{ld}}$ / $n_{\mathrm{st}}$ (doubleword) & 0.60 & 1.05 & 1 & 1 & 1 \\
$n_{\mathrm{ldsub}}$ / $n_{\mathrm{stsub}}$ (sub-doubleword) & 0.60 & 1.10 / 1.13 & 1 & 3 / 9 & 1 \\
$n_{\mathrm{br}}$ & 1.00 & 1.09 & 1 & 1 & 1 \\
$n_{\mathrm{jmp}}$ (jal, jalr) & 0.73 & 1.00 & 1 & 1 & 1 \\
$n_{\mathrm{sys}}$ (HALT, SPAWN, JOIN) & 1 & 1 & 1 & 1 & 1 \\
$m_{\mathrm{words}}$ (first-touched doublewords) & 0 & 0.12 & 0 & 0 & 0 \\
$m_{\mathrm{pages}}$ (first-touched 1\,KiB pages) & 0 & 0 & 1{,}130 & 0 & 0 \\
$P_{H,P}$ per compression (level: $+35$) & 72 & 174 & 100\,[J] & 200\,[J] & 8 \\
$\lambda_P$ per proof byte & 1/8 & 1/8 & 1/8 & 1/8 & 1/8 \\
\bottomrule
\end{tabular}
\caption{zkVM profiles of the recursion board [D4 §2a]. [J] marks a judgement. The three
non-ZisK Zbc entries assume a carry-less-multiply chip that SP1, RISC~Zero and Jolt do not have today.
The RISC~Zero and Jolt hash prices are guesses, to be replaced from published accelerator costs.}
\label{tab:cost-profiles}
\end{table}

The spread \emph{between} zkVMs on ALU-class operations is only 0.7--1.4 times. On memory footprint
(0 to 1{,}130 per KiB), sub-word stores (1 to 9), division (1.3 to 8) and hashes (72 to 200 per
compression, against 8 in v0.2) it is 10--100 times. These are exactly the counters v0.2 lacked.

\paragraph{\texttt{QUERY}: a virtual Merkle opening.} \texttt{QUERY} returns a whole declared leaf of
$\mathsf{leaf}_j$ $K$-elements of oracle $j$, which has length $\ell_j$. With
$\mathsf{depth}_j=\log_2(\ell_j/\mathsf{leaf}_j)$ (both powers of two) and
$b_j=\mathsf{leafbytes}_j=8\,\mathsf{leaf}_j$, a compiled verifier hashes the leaf, $\lceil b_j/64\rceil$
compressions, and one node per level. The exact price of one call under profile $P$ is therefore
\[
  \mathsf{depth}_j\cdot(P_{H,P}+35)+P_{H,P}\cdot\lceil b_j/64\rceil .
\]
The $+35$ is the measured non-hash cost of one level. The certificate does not carry
$\sum_i\lceil b_i/64\rceil$, so the contract charges the leaves of a run through the upper bound
\[
  \sum_i\Bigl\lceil\frac{b_i}{64}\Bigr\rceil\;\le\;\sum_i\Bigl(\frac{b_i}{64}+1\Bigr)
  \;=\;\frac{q^{\mathrm{bytes}}}{64}+q^{\mathrm{count}},
  \qquad\text{charged as } P_{H,P}\cdot\Bigl(\frac{q^{\mathrm{bytes}}}{64}+q^{\mathrm{count}}\Bigr),
\]
which uses only existing counters and over-charges by at most one compression per \texttt{QUERY}
(\lean{Profile.C}; the docstring gives the argument). Returning whole leaves is deliberate. With
single-element queries, a $2^4$-fold query would pay 16 paths. An amortised per-element price would be
gamed by declaring huge leaves and reading one element of each. A depth-0 round is a plain read,
priced by its leaf hash and its bytes. Every call is charged; caching is the verifier's own job.

\paragraph{The leaf-pricing hole closed in round~3 (A6-P1).} v0.2 charged the leaf as
$q^{\mathrm{bytes}}/8$, exactly the leaf hash at v0.2's $P_H=8$; v0.3 first rescaled $P_H$ per profile
for the levels and the absorption but left the leaf at that rate, 9--25 times below its hash cost.
Review A6 measured the consequences with the contract's own functions (\path{probes/a6-probes/Pricing.lean},
re-run in CI): moving a 1{,}024-word $\vk$ into the index oracle undid the A4-F6 fix, lowering the
$C+\mathit{FS}+\lceil\lambda B\rceil$ term by 7--45\%, and the charged score-optimal leaf was
$2^6$--$2^7$ against a true optimum of $2^4$. Under the corrected price the move costs exactly one
$P_{H,P}$ more than tape absorption (theorem \lean{vk\_in\_idx\_no\_gain}; RISC~Zero more, through
paging), the charged optimum is $\mathsf{leaf}=2^4$ in every profile, and a kernel-checked lemma gives
$\text{exact}\le\text{charged}\le\text{exact}+P_{H,P}$ for $\mathsf{leaf}\le2^{12}$ (row C25 has the
numbers). The latency board still charges the leaf one cycle per element (M-6).

\paragraph{Fiat--Shamir surcharge, per profile.} In the IT model challenges are free tape contents. A
compiled verifier absorbs the statement and every root, and squeezes every challenge. The contract
fixes the challenge menu: $E$-vectors and bounded indices, with no submitter-chosen types. That removes
the free \emph{derived} values a submitter type would let the harness compute, such as barycentric
weights, powers, inverses and $\mathrm{eq}$ tables [A1-M1]. What remains is priced statically from the
declared shapes:
\[
  FS_{P,k}=\mathrm{rounds}_k\,(2P_{H,P}+30)+n^{E}_k\,(0.75\,P_{H,P}+6)+4\,n^{\mathrm{idx}}_k
  +\Bigl\lceil\tfrac{(\mathsf{vkLen}+n_x)\cdot8}{64}\Bigr\rceil\,P_{H,P}.
\]
\begin{itemize}[nosep]
\item Per round: absorb a 32-byte root (one compression plus about 30 ALU operations) and squeeze one
  block.
\item Per $E$ element: 24 squeezed bytes, i.e.\ 0.75 compressions, plus loads. No reduction is needed,
  because every 24-byte string is an $E$ element.
\item Per index challenge: mask $b$ bits.
\item Per class: absorb the prefix $\vk_i\,\|\,x$, at most $256$ compressions at
  $\mathsf{vkLen}=n_x=2^{10}$. A4 found that v0.2 omitted this term [A4-F6].
\end{itemize}
At $P_H=8$ the first three terms reduce to v0.2's $46/12/4$.

\paragraph{Bytes and blow-up.} The byte count is
$B_k=\mathsf{tapeBytes}_k+q^{\mathrm{bytes}}_k+32\,q^{\mathrm{lev}}_k+32\cdot\#(\text{msg rounds})_k$.
It charges one digest per virtual level and one root per round, which closes the free-Merkle-cap
channel [A1-m3]. The charge $\lambda_P=1/8$ is a monotonicity device; small proofs are rewarded by the
proof-size board. RISC~Zero's paging now charges first-touched proof bytes through $m_{\mathrm{pages}}$,
about 1.1 cycles per byte. Rate is priced by
\[
  V_k=\sum_{j\ge1}\frac{\ell_j}{N_k+2^{k+4}}\quad(\idx\text{ excluded}),\qquad
  \text{factor }\max(1,V_k/16)^{1/4},\qquad V_k\le2^8 .
\]
Without it, every record would sit on the rate cap [A1-M4]. Threshold 16 makes SPARK at $\rho=1/2$
free. Exponent $1/4$ turns rate into a visible slope: exponent 0 puts the optimum at the cap and
exponents of $1/2$ or more force $\rho=1/2$.

\phantomsection\label{subsec:cost-score}%
\paragraph{The score.} For profile $P$ and class $k\in\{16,20,24\}$, with $c_k$ certified by
$\mathsf{AcceptingCostAtMost}(k,c_k)$:
\begin{align*}
  S_{P,k}&=\Bigl(w_P\cdot c_k+(P_{H,P}+35)\,q^{\mathrm{lev}}_k
           +P_{H,P}\bigl(q^{\mathrm{bytes}}_k/64+q^{\mathrm{count}}_k\bigr)+FS_{P,k}
           +\lceil\lambda_PB_k\rceil\Bigr)\cdot\max(1,V_k/16)^{1/4},\\
  \mathrm{Score}_{zk}&=\Bigl(\prod_{P\in\{A,B,C,D\}}\prod_{k}S_{P,k}\Bigr)^{1/12},\qquad
  \mathrm{Score}_{T2}\text{: the same with }S_{\mathrm{leaf}}+S_{\mathrm{fold}}+S_{\mathrm{dec}}/64 .
\end{align*}
Lower is better. Three properties make it hard to game.
\begin{itemize}[nosep]
\item \emph{Normalisation cancels.} Rankings under a geometric mean do not depend on per-profile
  normalisation constants, because those cancel in the ratio between two submissions.
\item \emph{Specialisation is visible and bounded.} Tanking one profile by a factor of 2 costs
  $2^{1/4}$, i.e.\ 19\%, overall, and the four per-profile scores are displayed.
\item \emph{The score is a theorem.} The scores are Lean functions of the certified data (\lean{scoreZk4}
  returns the fourth powers in $\mathbb Q$). The challenge renders the image bytes and every
  shape-derived input, and the certificate proves the claimed per-profile values. The external scorer
  only takes roots [A4-F2, F4].
\end{itemize}

\begin{decision}{COST-2}{Seasonal re-weighting without re-proof}
The profile family and its aggregation are announced. The table is published with the range each
entry may take, e.g.\ $n_{\mathrm{mul}}\in[1,1.5]$, $m_{\mathrm{pages}}\in[0,1500]$,
$P_H\in[60,250]$. Constants are refreshed per season, only from published zkVM cost tables. Every
record is re-scored from its certified vector; nothing is re-proved. If agents converge on one
season's constants, a per-season draw inside the ranges is held in reserve, which makes minimax over
the ranges the robust strategy. \emph{Rationale:} announcing the family prevents single-model
overfitting. Hiding exact constants buys little, because agents infer them from the leaderboard
within a season, and it costs transparency. [D4-6]
\end{decision}

%-----------------------------------------------------------------------------------------------
\subsection{The latency board: a deterministic timing model with fork-join harts}\label{subsec:cost-lat}

Native verification latency matters for the decider and for node operators. Its largest levers are
pipeline stalls, branch behaviour and parallelism, and per-instruction counts see none of them. The
latency board uses a model that is deterministic and therefore provable, and calibrates it against
hardware.

\paragraph{Scoreboard timing model.} The model is an in-order pipeline:
\begin{itemize}[nosep]
\item the state \lean{SB} is the next free issue slot (two per cycle) and a ready time per register;
\item an instruction issues in the first free slot at or after the cycle its sources are ready, and
  sets its destination ready after its kind's latency;
\item a taken branch or jump adds a bubble that flushes the rest of the cycle;
\item there are no caches.
\end{itemize}
The timer is a pure function of the existing trace: \code{Timer.lean} (22 code lines) plus the cost
classes, latencies and the \lean{tock} threading in \code{Machine.lean}. Latencies follow GCC's
SiFive-7 and SpacemiT-X60 scheduler models: ALU 1, load 3, MUL and clmul 3, a divide of 16, and a
branch penalty of 5. (The design asked for a non-pipelined divider; the compiled timer pipelines it,
see item~9 of I-COST-1.) The key property is compositionality, proved for the contract's dual-issue
scoreboard in the untrusted \code{Library/Timer.lean} (174 lines; \lean{Ok ops} says every latency
is at most $L_{\max}=16$, which holds for every operation the machine issues):
\begin{leancode}[label={\footnotesize Library/Timer.lean (statements)}]
theorem run_slack {sb : SB} (h : Bounded sb) (ops : List Op) :
    (run sb ops).time ≤ sb.time + time ops + Lmax       -- from any reachable entry state
theorem time_append (a b : List Op) (ha : Ok a) : time (a ++ b) ≤ time a + time b + Lmax
theorem time_blocks : ∀ (bs : List (List Op)), (∀ b ∈ bs, Ok b) →
    time bs.flatten ≤ (bs.map time).sum + Lmax * bs.length
theorem length_le_time (ops : List Op) : ops.length ≤ 2 * time ops   -- non-vacuity
\end{leancode}
A submitter proves per-block times from a fresh state and composes them with a constant slack per
block boundary. No reasoning about entry states is needed. Loops cost $n\cdot(\text{time of the
body}+L_{\max})$. The baseline proved its latency clause with a timer potential instead, at about 170
of its 1{,}920 lines, i.e.\ about 10\% on top of the counter certificate (\cref{subsec:base-proof}). The
deterministic model has precedent in time-predictable processors (PRET, Patmos), whose timing is part of
the ISA specification precisely so that worst-case analysis composes.

\paragraph{\texttt{QUERY} in the timer.} A \texttt{QUERY} stalls the pipeline for
$\mathtt{QUERY\_LEVEL\_CYCLES}\cdot\mathsf{depth}_j+\mathsf{leaf}_j$ cycles, with
$\mathtt{QUERY\_LEVEL\_CYCLES}=1{,}560$ per level (one BLAKE2s compression on the in-order model). The
leaf term, one cycle per element, is about 195 times below its hash cost; this latency-board half of
A6-P1 lives in the frozen machine timer and is a v0.4 item (M-6), since \lean{scoreLat3} sees only
$(W,S)$.

\paragraph{Fork-join harts.} Two syscalls add deterministic parallelism.
\begin{itemize}[nosep]
\item \texttt{SPAWN} ($t_0=3$) takes an entry point, an argument, the child's stack pointer, a write
  region $W_r$ and a read region $R$ (all 8-aligned; the child's stack must lie inside $W_r$).
  \texttt{JOIN} ($t_0=4$) takes the returned handle.
\item In Lean the semantics are \emph{sequentialised}: at \texttt{SPAWN} the child runs to completion
  at once, by a nested \lean{execute}, from zeroed registers with \texttt{a0} set to the argument. Its
  loads outside $R\cup W_r$ fault, and its stores and \texttt{QUERY} destinations outside $W_r$ fault.
  It must end with \texttt{HALT(1)}, otherwise the whole run rejects. Its stores become visible to the
  parent only at \texttt{JOIN}.
\item While a child is live, no other hart may store into its $R\cup W_r$ or load from its $W_r$; the
  regions of a new child must be accessible to its parent under these rules. Read regions of different
  children may overlap; write regions may not overlap any live child's regions. A violation faults, and
  a fault is a reject.
\end{itemize}
Under these rules every real interleaving reaches the same state at \texttt{JOIN} as the sequential
semantics (argued, not proved in Lean; A6-N5, L-10). A deployed runtime may therefore run children on
any cores in any order, and no concurrency enters the trusted model. A submitter needs no race-freedom
proof: a race is a fault, a faulting run rejects, and $\mathsf{Agrees}$ forces honest runs to be
race-free. Caps: at most $2^{16}$ spawns per
run, 64 live children per hart, nesting depth 8. Each \texttt{SPAWN} and each \texttt{JOIN} costs
$c_{\mathrm{spawn}}=c_{\mathrm{join}}=2{,}000$ pipeline cycles, about 1\,\textmu s for a pooled
runtime.

\paragraph{Work, span and the score.} Every hart carries its counter vector and its pipeline time.
\begin{itemize}[nosep]
\item \emph{Work.} A child's counters accumulate in the run's single accounting state, threaded
  through the nested \lean{execute}, so they are charged at \texttt{SPAWN} time. The recursion board's
  certificate keeps its shape: a zkVM proves total work $W$ and shards it, so span is irrelevant there.
  On that board a \texttt{SPAWN} and a \texttt{JOIN} count one $n_{\mathrm{sys}}$ each, at weight~1,
  which under-prices the context switch a deployed zkVM image needs (A6-H1, v0.4 item M-5).
\item \emph{Span.} A parent records $S_0=S_{\mathrm{parent}}$ at \texttt{SPAWN}. At \texttt{JOIN} it
  sets $S_{\mathrm{parent}}\leftarrow\max(S_{\mathrm{parent}},\ S_0+c_{\mathrm{spawn}}+S_c+c_{\mathrm{join}})$,
  where $S_c$ is the child's own span. This is the series-parallel critical path.
\item \emph{Score.} Brent's bound $T_p\le W_{\mathrm{pipe}}/p+S$ holds for any greedy scheduler, and
  work stealing achieves it in expectation up to the constant on $S$. The board scores
\[
  T_{16,k}=W_{\mathrm{pipe},k}/16+S_k,\qquad \mathrm{Score}_{\mathrm{lat}}=\Bigl(\prod_kT_{16,k}\Bigr)^{1/3},
\]
  certified by the optional clause \lean{AcceptingLatencyAtMost} on accepting runs. $T_{64}$ is
  displayed.
\end{itemize}
The board rewards query-parallel designs: openings, batched inversions, reduction trees. It penalises
long sequential sumcheck chains, but only on this board. In the trajectory model $S/W$ is 0.01--0.24,
so $T_{16}/W$ is 0.07--0.30 (\cref{subsec:bench-trajectory}). SIMD (RVV) is deferred. No zkVM
implements it, its semantics would multiply the size of the machine file, and harts already capture
the data parallelism.

\paragraph{Wall time: calibration and display.} Each season the organisers run the calibration kernels
on a reference RISC-V board with Zbc (SpacemiT K1 class; its Zbc GHASH runs 19 times faster than
without) and fit the latency table by least squares. Pinned frequency, isolated cores and at least 20
interleaved runs per measurement keep the noise at $\pm2$--$5\%$. Measured times are shown next to the
certified bound. A submission whose board time deviates from its $T_{16}$ by more than a factor of 2 is
flagged for review. Wall time is never the score, for three reasons. Records are decided by sub-percent
margins, so \enquote{submit many, keep the lucky run} would pay. A measurement is not a theorem. And a
measured score would put candidate code back on the score's trust path.

%-----------------------------------------------------------------------------------------------
\subsection{Admission caps}\label{subsec:cost-caps}

Every cap that a score or the harness relies on is a decidable certificate clause rendered into the
statement [A2-m3]. \Cref{tab:cost-caps} lists the caps. Lookup-table caps stop free precomputation,
prover caps stop the rate race, protocol caps keep the security caps from being made vacuous, and
hart caps stop spawn storms.

\begingroup\small
\renewcommand{\arraystretch}{1.1}
\begin{longtable}{@{}>{\raggedright\arraybackslash}p{0.25\textwidth}>{\raggedright\arraybackslash}p{0.13\textwidth}>{\raggedright\arraybackslash}p{0.55\textwidth}@{}}
\caption{Admission caps and the failure each answers.}\label{tab:cost-caps}\\
\toprule
Cap & Value & Purpose (attack answered) \\
\midrule
\endfirsthead
\toprule
Cap & Value & Purpose (attack answered) \\
\midrule
\endhead
\multicolumn{3}{@{}l}{\emph{Lookup-table caps}}\\
code + data of each image & $\le2^{20}$\,B & an unbounded image is a free lookup table (ots.golf's post-launch cap). Code is now also priced, through $m_{\mathrm{pages}}$ and the timing model \\
verification key & $\le2^{10}$ $K$-elem. & a large key is a per-circuit table; the index lives in the priced oracle $\idx(i)$ \\
accumulator plain part; oracles & $\le2^{12}$; $\le16$ & concatenation batchers; a free table carried between folds (\cref{subsec:rec-shape}) \\
tape & $\le2^{20}$\,B & bounds the challenge channel (about 43k $E$ elements) \\
public input & $n_x=2^{10}$ & with $n_x=2^{16}$, ingesting $x$ was 27--50\% of every score [A1-M2] \\
\multicolumn{3}{@{}l}{\emph{Prover-cost caps}}\\
oracle length & $\ell_j\le2^{k+10}$ & sanity bound; with $\mathsf{leaf}_j\mid\ell_j$ it fixes $\mathsf{depth}_j$ [A1-m10] \\
blow-up & $V_k\le2^8$ & hard cap behind the soft factor \\
\multicolumn{3}{@{}l}{\emph{Protocol caps}}\\
rounds & $\le2^{14}$ & without it the per-round error cap bounds nothing (\lean{hitRed\_secure}) [A2-C3]; $2^{10}$ would ban GKR-based SPARK [A1-M6] \\
queries & $Q_k\le2^{16}$ & bounds data-dependent query loops \\
per-round error & $\le2^{-128}$ & a subfield challenge cannot meet it (\cref{sec:security}) \\
index oracle & $\mathit{idxLen}\le2^{k+8}$ $K$-elem. & bounds the threshold $t^*<2^{43}$ of the compiled bound (\cref{subsec:sec-cap}) [A5-M4] \\
decider & none beyond $F$ & no separate cap: linear (Arc/WARP-style) deciders are admissible and enter T2 at weight $1/64$ [A5-M8] \\
fuel & $F=2^{36}$ per hart & termination by construction; bounds DoS at $2^{36}$ steps without harts and $(\#\mathrm{spawns}+1)\cdot2^{36}$ with harts (H12) \\
\multicolumn{3}{@{}l}{\emph{Hart caps (latency board)}}\\
spawns; live harts; depth & $2^{16}$; 64; 8 & spawn storms; plus $c_{\mathrm{spawn}}=c_{\mathrm{join}}=2{,}000$ on both $W$ and $S$ \\
\bottomrule
\end{longtable}
\endgroup

\paragraph{Prover-cost constraints.} Completeness needs only \emph{some} honest prover, a Lean
\lean{Strategy}. Prover cost still has to be bounded, or the rate goes to zero. The instruments are:
\begin{itemize}[nosep]
\item the blow-up factor and caps above (static and decidable);
\item the computability gate on the badge-path fields \lean{VK}, \lean{IDX}, prover and extractor,
  run on the elaborated environment [A4-F3, A7-M2]; it only guarantees that compiled code exists,
  since a brute-force table is Lean-computable (C-7);
\item badge-path budgets, which are what bound feasibility: 1\,h at $2^{20}$ and 8\,h at $2^{24}$ on the
  reference machine, under a RAM budget (\cref{subsec:bench-badge}; not yet implemented, H-5).
\end{itemize}
The IOP track has true randomness, so it cannot reward grinding (D2-G6; the ROM track, Q-4, restores it).

\paragraph{Implementation cost of the score.} The brief asks for a cost function that is easy to
implement. The scorer is small and has no second implementation to drift: the cost clauses, the five
profiles, $C_P$, the shape summary with its FS, byte and blow-up terms, \lean{scoreZk4},
\lean{scoreLat3} and the admission predicates are 73 code lines of \code{Cost.lean}
(\cref{subsec:app-cost}), and the certificate proves that the claimed scores are their values, so
the external scorer only takes roots. A claim carries the 18-counter vector per class (and $W,S$ per
class for a latency record) next to the image; the shape inputs are rendered from the declared
protocol. Re-weighting needs no re-proof. A minimal launch can omit
the latency board and harts (both optional clauses: \lean{Latency1}, \lean{img.harts}), and the
v0.2 weights are a display column only. The proof burden is the 18-counter bound, which the
baseline's 1{,}920 lines show to be routine for a straight-line verifier.

%-----------------------------------------------------------------------------------------------
\subsection{Field-arithmetic costs and their incentives}\label{subsec:cost-field}

\begin{table}[t]
\centering\footnotesize
\setlength{\tabcolsep}{3pt}
\begin{tabular}{@{}lrrrrrrrr@{}}
\toprule
Kernel (units per operation) & uni & zisk & sp1z & jolt & r0 & sp1 & inord & ooo \\
\midrule
$K$ mult.\ (shift/XOR reduction) & 35 & 26.8 & 34.0 & 29.0 & 58.2 & 381 & 20.1 & 12.1 \\
$E$ mult.\ (Karatsuba) & \textbf{139} & 98.8 & 123.3 & 103 & 212 & 2{,}207 & 65.2 & 25.9 \\
$E$ mult.\ (schoolbook) & 157 & 98.7 & 124.7 & 103 & 206 & 3{,}250 & 74.2 & 25.9 \\
$K\times E$ MAC, per term & \textbf{38.1} & 17.8 & 24.1 & 20.1 & 63.4 & 1{,}066 & 17.8 & 5.2 \\
sumcheck round, deg.\ 3, over $E$ & 2{,}522 & 1{,}754 & 2{,}210 & 1{,}874 & 2{,}824 & 39{,}712 & 1{,}255 & 485 \\
$\mathrm{eq}(r,u)$, 24 variables & 3{,}303 & 2{,}365 & 2{,}912 & 2{,}439 & 3{,}815 & 52{,}915 & 1{,}536 & 624 \\
$K$ inverse, Itoh--Tsujii & 2{,}053 & 1{,}583 & 1{,}951 & 1{,}627 & 2{,}762 & 26{,}605 & 1{,}241 & 944 \\
$K$ inverse, Fermat & 3{,}506 & 2{,}694 & 3{,}322 & 2{,}756 & 4{,}148 & 46{,}727 & 1{,}855 & 740 \\
BLAKE2s compression (software) & 2{,}662 & 2{,}261 & 3{,}004 & 3{,}744 & 3{,}499 & 3{,}004 & 1{,}560 & 709 \\
Merkle level, model price & 43 & 107 & 209 & 235 & 135 & 209 & 1{,}595 & 744 \\
trivial baseline $2^{20}$ (M) & 329 & 229 & 301 & 285 & 497 & 3{,}942 & 1{,}050 & 214 \\
\bottomrule
\end{tabular}
\caption{Field kernels under the cost models [D4 RESULT §1.1, §1.3, §2]. \texttt{sp1} has no Zbc and
emulates carry-less products by tables; \texttt{sp1z} assumes a Zbc chip priced like MUL. Native rows
(\texttt{inord}, \texttt{ooo}) are literature-latency models, not silicon measurements.}
\label{tab:cost-field}
\end{table}

\Cref{tab:cost-field} prices the kernels that dominate succinct verifiers. Five consequences follow.
\begin{enumerate}[nosep]
\item \emph{Zbc is a deployment assumption.} Only ZisK executes \texttt{clmul} natively today. Without
  a carry-less chip, every binary-field kernel is 11--28 times dearer (\texttt{sp1}), and FRI-shaped
  designs with the fewest field operations rank first. The challenge text states this assumption, and
  the profiles do not hedge it.
\item \emph{The Karatsuba lever disappears.} Karatsuba beats schoolbook by 13\% under \texttt{uni} and
  ties under \texttt{zisk}, \texttt{jolt} and \texttt{ooo}. Instruction-count tricks in field code are
  not the levers v0.2 suggested. The score tracks architecture more and assembly less.
\item \emph{Inner products are cheap.} Carry-less sums never overflow, so they accumulate by XOR with
  one reduction at the end: 18--24 units per $K\times E$ term on the ZisK, SP1 and Jolt profiles. Hint-checked inversions
  (one multiplication plus 24 bytes) remain right everywhere. The Fermat-versus-Itoh--Tsujii choice
  flips only under out-of-order execution.
\item \emph{The hash-to-field ratio changes.} A Merkle level is now worth 0.6--2.3 $E$ multiplications on
  the zkVM profiles, against 0.3 in v0.2, and about 25 on native cores. Query-heavy designs pay
  accordingly.
\item \emph{Tables still cannot buy multiplication back, given Zbc.} The 16\,KiB squaring table loses
  1.8--2.5 times under every model with Zbc and wins 6 times only without it.
\end{enumerate}

%-----------------------------------------------------------------------------------------------
\subsection{Sensitivity analysis}\label{subsec:cost-sensitivity}

\begingroup\small
\renewcommand{\arraystretch}{1.1}
\begin{longtable}{@{}>{\raggedright\arraybackslash}p{0.24\textwidth}>{\raggedright\arraybackslash}p{0.46\textwidth}>{\raggedright\arraybackslash}p{0.23\textwidth}@{}}
\caption{Sensitivity of the rankings to models and constants [D4 RESULT §2, §5; A1].}\label{tab:cost-sensitivity}\\
\toprule
Model or constant & Effect & Rank flips? \\
\midrule
\endfirsthead
\toprule
Model or constant & Effect & Rank flips? \\
\midrule
\endhead
choice of zkVM profile & absolute levels 0.9 (ZisK) to 2.0 (RISC~Zero) times v0.2; top three architectures stable & FRI-shaped designs drop 1--2 ranks under hash-priced profiles \\
no Zbc chip (\texttt{sp1}) & binary kernels $\times$11--28 & yes: FRI-shaped design rank 1 \\
latency board vs recursion board & GKR-style SPARK (588 sequential rounds) rank 4$\to$7; FRI-shaped 5$\to$4; UD 6$\to$5 & yes, by design \\
pipeline parameters & single issue moves field kernels from 0.5 to 0.8 of v0.2; clmul latency 2 vs 5 gives $\pm20\%$; DRAM 200 doubles random-load costs & jump-table-vs-branch flip reverses \\
blow-up exponent $0,\tfrac14,\tfrac12,1$ & optimal $\rho$: cap, slope, $1/2$, $1/2$ [A1-M4] & yes: chooses the rate \\
$\lambda\in[0,8]$ & bytes are 0.3--0.6\% of succinct scores [A1-m1] & none \\
$/64$ vs $/1024$ (T2) & a linear decider beats a succinct one only if the succinct fold is $\ge30\%$ dearer [A1-m5] & at the margin \\
classes $\{16,20,24\}$ & crossover with the trivial verifier at $k\approx13$; geomean $\approx$ the $2^{20}$ class [A1-m2] & possible for designs with steep growth in $k$ \\
proximity regime & UD $\times1.7$ over JB; capacity conjecture (banned) $\times0.53$ & by track \\
\bottomrule
\end{longtable}
\endgroup

%-----------------------------------------------------------------------------------------------
\subsection{What the cost function fails to capture}\label{subsec:cost-fails}

Three gaps are structural to the IOP track and are not addressed by v0.3.
\begin{enumerate}[nosep]
\item \emph{BCS engineering is virtual.} Pruned multi-openings, Merkle caps and arity, digest
  truncation and grinding are not design axes. They decide roughly 1.5--2 times of real hash cost. The
  ROM track restores them, with \texttt{HASH} priced per profile like the virtual level.
\item \emph{Prover cost} is bounded only by static oracle lengths and unscored demonstration budgets.
\item \emph{Recursion cost is in zkVM work, not in constraints of an arithmetised verifier} (T3,
  \cref{subsec:rec-t3}). A leanISA step is worth 1--139 RV64IM+Zbc cycles depending on the opcode, so
  field-native VMs are not directly comparable.
\end{enumerate}
The modelling gaps that v0.3 narrows but does not close are listed in the box.

\begin{issue}{I-COST-1}{What remains unrealistic in v0.3}
\begin{enumerate}[nosep]
\item \emph{No cache, TLB or DRAM model.} Random-access verifiers are under-priced by up to 7 times
  (kernels) and 3 times (the trivial baseline at $2^{20}$) on in-order cores. First-touch counts cover
  footprint, not locality. This is immaterial for succinct verifiers with a working set under 1\,MiB,
  but a hole for table-driven designs.
\item \emph{No instruction-level parallelism beyond dual issue.} The 1.5--2 times of native performance
  that independent chains buy is invisible.
\item \emph{One latency table} stands for \enquote{an in-order core}. Constants vary 2--3 times across
  cores.
\item \emph{Zbc native in every zkVM profile}; true only of ZisK.
\item \emph{Two hash prices are guesses} (RISC~Zero, Jolt), and SP1's fetch and decode areas are
  assumed per executed instruction.
\item \emph{Protocol counts} in the trajectory carry the $\pm2\times$ model uncertainty.
\item \emph{Latency-board leaf hashing} (A6-P1, second half): the timer charges a \texttt{QUERY}
  leaf one cycle per element instead of $1{,}560$ per 64 bytes.
\item \emph{\texttt{SPAWN}/\texttt{JOIN} on the recursion board} (A6-H1): one $n_{\mathrm{sys}}$
  each, against a context switch of about 62 memory operations in a deployed zkVM image.
\item \emph{Timer structural hazards} (A6-T1): two load/store slots per cycle, a fully pipelined
  16-cycle divider and no WAW stall, so memory-bound loops (Merkle path copies, table-driven field
  code) look about 2 times faster than on an in-order core with one memory port. Proposal: one
  memory port and an unpipelined divider (about 6 trusted lines), or calibration on the board.
\end{enumerate}
Items 7--9 are scheduled for v0.4 (DESIGN §15.4) together with the hart fuel rule of
\cref{subsec:thr-catalogue} (row H12).
Actions: verify the SP1 and RISC~Zero facts and add OpenVM and Airbender profiles. Fit the latency
table on a SpacemiT K1 and a P550. Investigate whether a 64-byte-line first-touch counter is provable
for query-driven access. Measure the FS constants on a BLAKE2s transcript. Re-tune the threshold 16
and the exponent $1/4$ against the first real submissions (G-D3-5, D4 §6).
\end{issue}
